\documentclass[12pt,a4paper]{article}

%% ---- Packages ----
\usepackage[margin=2.5cm]{geometry}
\usepackage{graphicx}
\usepackage{booktabs}
\usepackage{longtable}
\usepackage{array}
\usepackage{xcolor}
\usepackage{listings}
\usepackage{hyperref}
\usepackage{amsmath}
\usepackage{amssymb}
\usepackage{enumitem}
\usepackage{titlesec}
\usepackage{fancyhdr}
\usepackage{microtype}
\usepackage{parskip}
\usepackage{mdframed}
\usepackage{fontenc}
\usepackage{inputenc}

%% ---- Colours ----
\definecolor{noteblue}{RGB}{30,100,180}
\definecolor{warnorange}{RGB}{200,100,20}
\definecolor{codegrey}{RGB}{240,240,245}
\definecolor{linkblue}{RGB}{10,80,160}

%% ---- Hyperref ----
\hypersetup{
  colorlinks = true,
  linkcolor  = linkblue,
  urlcolor   = linkblue,
  citecolor  = linkblue,
  pdftitle   = {AGN Variability Archive -- Pipeline Manual},
  pdfauthor  = {J.V. Hernandez Santisteban},
}

%% ---- Code listings ----
\lstdefinestyle{shell}{
  backgroundcolor=\color{codegrey},
  basicstyle=\ttfamily\small,
  breaklines=true,
  frame=single,
  rulecolor=\color{gray!40},
  numbers=none,
}
\lstdefinestyle{python}{
  backgroundcolor=\color{codegrey},
  basicstyle=\ttfamily\small,
  breaklines=true,
  frame=single,
  rulecolor=\color{gray!40},
  language=Python,
  keywordstyle=\color{noteblue}\bfseries,
  commentstyle=\color{gray},
  stringstyle=\color{warnorange},
  numbers=left,
  numberstyle=\tiny\color{gray},
  numbersep=5pt,
}

%% ---- Headers / footers ----
\pagestyle{fancy}
\fancyhf{}
\fancyhead[L]{\textsc{AGN Variability Archive}}
\fancyhead[R]{\textsc{Pipeline Manual}}
\fancyfoot[C]{\thepage}
\renewcommand{\headrulewidth}{0.4pt}

%% ---- Section style ----
\titleformat{\section}{\large\bfseries\color{noteblue}}{}{0em}{}[\titlerule]
\titleformat{\subsection}{\normalsize\bfseries}{}{0em}{}

%% ---- Note / warning boxes ----
\newmdenv[
  linecolor=noteblue,
  backgroundcolor=noteblue!5,
  roundcorner=4pt,
  skipabove=6pt,skipbelow=6pt,
  innerleftmargin=8pt,innerrightmargin=8pt,
]{notebox}
\newmdenv[
  linecolor=warnorange,
  backgroundcolor=warnorange!5,
  roundcorner=4pt,
  skipabove=6pt,skipbelow=6pt,
  innerleftmargin=8pt,innerrightmargin=8pt,
]{warnbox}

%%============================================================
\begin{document}
%%============================================================

\begin{titlepage}
\centering
\vspace*{3cm}
{\Huge\bfseries AGN Variability Archive}\\[0.5em]
{\Large\bfseries Photometric Pipeline Manual}\\[2em]
{\large Version 1.0 --- November 2026}\\[1em]
{\large J.V.\ Hernandez Santisteban}\\[0.5em]
{\normalsize School of Physics \& Astronomy, University of St Andrews}\\[3em]

\begin{abstract}
\noindent
This manual documents the photometric data-reduction pipeline of the AGN Variability Archive
(AVA), a multi-band monitoring programme conducted with the Las Cumbres Observatory (LCO)
global telescope network.  The pipeline has been in production since 2021 and processes light
curves for approximately 70 AGN supporting reverberation mapping campaigns.  Six stages are
described in detail: orchestration, initialisation, data download, source extraction with
SExtractor, photometric calibration against the APASS catalogue, and inter-calibration across
telescopes with the CALI algorithm.  Outputs include calibrated flux light curves in mJy,
diagnostic plots, and a merged light curve.  The manual also catalogues known limitations and
maps each stage to its replacement in AVA~v2.
\end{abstract}

\vfill
{\small\textit{Internal technical document.  Not for distribution.}}
\end{titlepage}

\tableofcontents
\clearpage

%%============================================================
\section{Overview}
\label{sec:overview}
%%============================================================

The AGN Variability Archive (AVA) photometric pipeline processes multi-band imaging light
curves from the Las Cumbres Observatory global telescope network for approximately 70 active
galactic nuclei.  The overarching scientific goal is reverberation mapping: measuring the
time lag between AGN accretion-disc continuum variability and the response in the broad
emission-line region, which yields black-hole mass estimates.

\subsection{Instruments and telescopes}

Photometry is obtained with two camera families:
\begin{itemize}
  \item \textbf{Sinistro} (fa-series detectors, 1\,m telescopes) at sites coj, ogg, lsc,
        cpt, elp, tfn.
  \item \textbf{Spectral} (fs-series detectors, 2\,m Faulkes telescopes) at ogg and coj.
\end{itemize}
Standard LCO broadband filters are used: $u'g'r'i'z'$ (SDSS system), plus $B$, $V$, $R$,
$I$ (Johnson--Cousins) for some targets.  FLOYDS spectrographs on the 2\,m telescopes
provide simultaneous spectroscopic monitoring handled by a separate pipeline
(see the FLOYDS RM report).  Some targets also have Liverpool Telescope (LT/IO:O)
photometry ingested through a parallel processing branch.

\subsection{Pipeline outputs}

For each target on each observing night the pipeline produces:
\begin{enumerate}
  \item Calibrated flux light curves (mJy, all filters), appended to \texttt{lco\_latest.pkl}.
  \item Publication-quality and diagnostic plots (\textsc{pdf} + \textsc{png}) copied to the
        AVA website.
  \item An inter-calibrated merged light curve correcting telescope-to-telescope flux offsets.
  \item Standard-star catalogues (\texttt{lco\_latest\_stan.pkl}) for internal photometric
        quality monitoring.
\end{enumerate}

\subsection{Data flow}

\begin{figure}[h!]
\centering
\begin{minipage}{0.85\linewidth}
\begin{lstlisting}[style=shell]
target_list.txt
      |
      v
ava_update.py  <-- flags: -o -d -e -c -p -m -s
  (one shell script per target, run with nohup)
      |
      +-> ava_initialise.py    set up directories, fetch APASS catalogue
      +-> download_data.py     pull BANZAI FITS from LCO archive API
      +-> ava_sextractor.py    SExtractor aperture photometry  ->  _cat.fits
      +-> ava_calibrate.py     match APASS, fit zeropoint  ->  lco_latest.pkl
      +-> ava_webplot.py       convert to mJy, plot all filters + Swift
      `-> ava_merge.py         CALI inter-calibration across telescopes
\end{lstlisting}
\end{minipage}
\caption{End-to-end data flow.  Each stage is an independent Python script
  invoked by a generated shell script; state passes via pickle files and FITS
  catalogues on disk.}
\label{fig:flow}
\end{figure}

\subsection{Technology stack}
Python~3, \texttt{numpy}/\texttt{scipy}, \texttt{astropy}, \texttt{pandas},
\texttt{matplotlib}, \texttt{lmfit}, SExtractor (external binary),
\texttt{pycali\_jvhs} (custom MCMC inter-calibration wrapper).

%%============================================================
\section{Stage 0 --- Orchestration: \texttt{ava\_update.py}}
\label{sec:orchestrator}
%%============================================================

\texttt{ava\_update.py} is the master driver.  It reads a plain-text target list, loops over
every target, writes one shell script per target that chains the requested pipeline stages,
and launches all scripts simultaneously using \texttt{nohup ... \&}.

\subsection{Target list}

A \textsc{csv} file (default \texttt{target\_list.txt}) with at minimum the columns:
\begin{itemize}[noitemsep]
  \item \texttt{targets} --- the LCO target-name string used in archive queries.
  \item \texttt{directory} --- the per-target folder slug, e.g.\ \texttt{Fairall\_9}.
  \item \texttt{LCO} / \texttt{LT} --- boolean flags enabling the LCO and LT branches.
\end{itemize}

\subsection{Command-line flags}

\begin{center}
\begin{tabular}{cl}
\toprule
Flag & Stage \\
\midrule
\texttt{-o} & Initialise directories (\texttt{ava\_initialise.py}) \\
\texttt{-d} & Download new data from LCO archive \\
\texttt{-e} & Source extraction with SExtractor \\
\texttt{-c} & Photometric calibration \\
\texttt{-p} & Web plots (\texttt{ava\_webplot.py}) \\
\texttt{-m} & Inter-calibration with CALI \\
\texttt{-s} & Save CALI merged output \\
\texttt{-i} & Colour image generation \\
\texttt{-r} & Remove \texttt{lco\_latest*} scratch files \\
\bottomrule
\end{tabular}
\end{center}

The \texttt{-project} argument selects the LCO key-project semester: \texttt{KEY2014A-002},
\texttt{KEY2018B-001}, \texttt{KEY2020B-006}, or \texttt{latest} (auto-detects the most
recent epoch folder present on disk and sets that night as the start date).

\subsection{Script generation}

For each target the driver opens \texttt{1script2rule\textit{dir}.sh}, writes the requested
stage commands as lines, closes the file, makes it executable, and launches it with
\texttt{nohup}.  Log output goes to \texttt{lco\_\textit{dir}\_logall.out}.

\begin{warnbox}
\textbf{Migration target.}  All scripts contain absolute paths hardcoded to
\texttt{/data/jvhs1/agn/} and \texttt{/data3/jvhs1/agn/}.  These must be replaced by
paths read from \texttt{config/ava\_config.yaml} in AVA~v2.
\end{warnbox}

%%============================================================
\section{Stage 1 --- Initialisation: \texttt{ava\_initialise.py}}
\label{sec:init}
%%============================================================

Run once per new target (\texttt{-o}).  Creates the directory tree under
\texttt{master\_dir} and fetches reference catalogues.

\subsection{Directory structure}
\begin{lstlisting}[style=shell]
<master_dir>/<obj>/          # working directory
<master_dir>/<obj>/UVOT/     # Swift/UVOT data placeholder
<master_dir>/raw_data/<obj>/ # raw FITS downloads
\end{lstlisting}

\subsection{APASS catalogue download}

The script uses the \texttt{mechanize} Python browser library to query the APASS
(AAVSO Photometric All-Sky Survey) web form for a cone search centred on the target's
RA/Dec.  The resulting CSV (\texttt{apass\_\textit{obj}.csv}) is the primary photometric
reference catalogue used in calibration.

\begin{warnbox}
\textbf{Migration target.}  This relies on a brittle web scrape that breaks on UI
changes.  AVA~v2 should use \texttt{astroquery.vizier.query\_region()} on APASS\,DR9,
or the Pan-STARRS PS1 catalogue via the LCO archive tools.
\end{warnbox}

%%============================================================
\section{Stage 2 --- Download: \texttt{download\_data.py}}
\label{sec:download}
%%============================================================

Queries the LCO science archive API (\texttt{archive-api.lco.global/frames/}) using an
authentication token for BANZAI-reduced frames (\texttt{rlevel=91}) and downloads them
to the target's raw data directory, organised by night (\texttt{YYYYMMDD/}).

\subsection{Date range logic}

When \texttt{-project latest} is set, the script inspects the most recent epoch folder
present on disk and uses that night as the start date, so only new data is fetched.  For
older semesters the date range is hardcoded (e.g.\ KEY2018B-001: 2018-06-01 to
2020-07-01).

\begin{notebox}
\textbf{AVA v2.}  The download step is replaced entirely by \texttt{lcodb sync} +
\texttt{lcodb daily}, which maintains a local metadata database without bulk-downloading
FITS files.  Images are fetched on demand at reduction time via the archive frame URL.
\end{notebox}

%%============================================================
\section{Stage 3 --- Source Extraction: \texttt{ava\_sextractor.py}}
\label{sec:sextractor}
%%============================================================

Runs SExtractor on every BANZAI-reduced FITS frame (\texttt{*e91.fits.fz}) found under
the target's raw data directory.

\subsection{Input preparation}

BANZAI frames are stored as compressed multi-extension FITS (\texttt{.fz}).  The script
opens each file with \texttt{astropy.io.fits}, extracts the \texttt{SCI} extension, writes
it to a temporary single-extension file (\texttt{temp\_sextractor\_file.fits}), and passes
that to SExtractor.  This step is necessary because SExtractor cannot handle
multi-extension FITS natively.

\subsection{SExtractor configuration}

Fixed configuration files (\texttt{default.sex}, \texttt{default.param},
\texttt{default.conv}, \texttt{default.nnw}, \texttt{default.psfex}) are stored in a
template directory and copied to each target's working folder at initialisation.  The key
output is a source catalogue in FITS binary-table format (\texttt{\_cat.fits}) written to
\texttt{sextractor\_ava/}.

\subsection{Apertures}

SExtractor is run with 12 aperture sizes simultaneously (\texttt{MAG\_APER},
\texttt{MAGERR\_APER}), producing a 12-element array per source.  The science aperture
index used in calibration depends on telescope and epoch (see Section~\ref{sec:aper}).

\subsection{Skip logic}

If a \texttt{\_cat.fits} file already exists and is non-empty for a given frame,
SExtractor is not re-run, making the stage incremental.

%%============================================================
\section{Stage 4 --- Photometric Calibration: \texttt{ava\_calibrate.py}}
\label{sec:calibrate}
%%============================================================

The core science stage.  Reads every BANZAI FITS frame and its SExtractor catalogue,
performs absolute aperture photometry on the AGN and on APASS standard stars, fits a
photometric zeropoint, and appends the result to \texttt{lco\_latest.pkl}.

\subsection{Incremental processing}

If \texttt{lco\_latest.pkl} already exists, only frames whose filename (\texttt{flname})
does not appear in it are processed.  The \texttt{-f} flag forces full reprocessing.

\subsection{Coordinate matching}

The AGN position is read from the target database.  SExtractor source positions are matched
using \texttt{SkyCoord.match\_to\_catalog\_sky()}.  Frames where the nearest source is more
than 3\,arcsec from the AGN position are skipped.

\subsection{Extinction and colour correction}

For each frame, site-specific extinction coefficients are read from
\texttt{lco\_ext\_coefficients.csv} (columns: \texttt{Filter}, \texttt{siteid},
\texttt{dome}, \texttt{K}, \texttt{C\_t}, \texttt{Color1}, \texttt{Color2}).
The corrected instrumental magnitude is:
\begin{equation}
  m_{\rm inst,corr} = m_{\rm inst} + 2.5\log_{10}(t_{\rm exp}) - K\,X - C_t\,c
  \label{eq:ext}
\end{equation}
where $X$ is airmass (minus 1), $t_{\rm exp}$ is exposure time in seconds, and $c$ is the
star's colour from APASS.  Both the AGN and each standard star receive this correction
before zeropoint fitting.


\subsection{Aperture correction}


\subsection{Aperture index resolution}
\label{sec:aper-index}

The SExtractor aperture series changed in 2025: two smaller diameters (3\,px
and 5\,px) were prepended to the original 10-aperture series, shifting every
old index by $+2$.  Using a fixed integer index across epochs therefore
samples \emph{different physical apertures} in the two epochs, corrupting the
light curves silently.

The correct solution, implemented in \texttt{ava.photometry.apertures}, is to
compute the index at reduction time from:
\begin{enumerate}
  \item the \emph{target aperture radius} (a configuration parameter,
        default  = 5$\,arcsec, i.e.\ 10\,arcsec diameter), and
  \item the \emph{pixel scale} read from the FITS header \texttt{PIXSCALE} key.
\end{enumerate}

\begin{lstlisting}[style=python]
from ava.photometry.apertures import science_aperture_index

# pixscale from FITS header; mjd from MJD-OBS keyword
idx = science_aperture_index(pixscale=0.389, mjd=61000.0)
# or from a full header dict:
from ava.photometry.apertures import science_aperture_index_from_header
idx = science_aperture_index_from_header(header_dict)
\end{lstlisting}

Tables~\ref{tab:aper-new} and \ref{tab:aper-old} list the physical aperture
diameters for every index and camera family.

\begin{table}[h!]
\centering
\small
\begin{tabular}{rr rrr}
\toprule
 & & \multicolumn{3}{c}{Physical diameter (arcsec)} \
\cmidrule(lr){3-5}
Index & Diam (px) & Sinistro 1\,m & Spectral 2\,m & MuSCAT3 2\,m \
      &           & (0.389\,\arcsec/px) & (0.304\,\arcsec/px) & (0.270\,\arcsec/px) \
\midrule
  0 &  3 &  1.17 &  0.91 &  0.81 \
  1 &  5 &  1.95 &  1.52 &  1.35 \
  2 &  8 &  3.11 &  2.43 &  2.16 \
  3 & 11 &  4.28 &  3.34 &  2.97 \
  4 & 13 &  5.06 &  3.95 &  3.51 \
  5 & 15 &  5.83 &  4.56 &  4.05 \
  6 & 18 &  7.00 &  5.47 &  4.86 \
  7 & 21 &  8.17 &  6.38 &  5.67 \
  8 & 24 &  9.34 &  7.30 &  6.48 \
  9 & 27 & 10.50 &  8.21 &  7.29 \
 10 & 30 & 11.67 &  9.12 &  8.10 \
 11 & 33 & 12.84 & 10.03 &  8.91 \
\bottomrule
\end{tabular}
\caption{New aperture series (2025+, 12 apertures).
  \texttt{PHOT\_APERTURES = 3,5,8,11,13,15,18,21,24,27,30,33} in
  \texttt{default.sex}.  Rows in bold indicate the index selected for a
  5\,arcsec-radius (10\,arcsec-diameter) science aperture on each camera.}
\label{tab:aper-new}
\end{table}

\begin{table}[h!]
\centering
\small
\begin{tabular}{rr rrr}
\toprule
 & & \multicolumn{3}{c}{Physical diameter (arcsec)} \
\cmidrule(lr){3-5}
Index & Diam (px) & Sinistro 1\,m & Spectral 2\,m & MuSCAT3 2\,m \
      &           & (0.389\,\arcsec/px) & (0.304\,\arcsec/px) & (0.270\,\arcsec/px) \
\midrule
  0 &  8 &  3.11 &  2.43 &  2.16 \
  1 & 11 &  4.28 &  3.34 &  2.97 \
  2 & 13 &  5.06 &  3.95 &  3.51 \
  3 & 15 &  5.83 &  4.56 &  4.05 \
  4 & 18 &  7.00 &  5.47 &  4.86 \
  5 & 21 &  8.17 &  6.38 &  5.67 \
  6 & 24 &  9.34 &  7.30 &  6.48 \
  7 & 27 & 10.50 &  8.21 &  7.29 \
  8 & 30 & 11.67 &  9.12 &  8.10 \
  9 & 33 & 12.84 & 10.03 &  8.91 \
\bottomrule
\end{tabular}
\caption{Old aperture series (pre-2025, 10 apertures).
  \texttt{PHOT\_APERTURES = 8,11,13,15,18,21,24,27,30,33} in
  \texttt{default.sex}.  Note: index~7 here (27\,px, 10.5\,arcsec on 1\,m)
  maps to index~9 in the new series.}
\label{tab:aper-old}
\end{table}

\begin{notebox}
\textbf{Science aperture for re-ingestion.}  For the half-million existing
frames, re-compute the science index via
\texttt{science\_aperture\_index(pixscale, mjd)} rather than storing a fixed
integer.  The function automatically selects the old (pre-MJD\,60909) or new
(MJD\,$\geq$\,60909) series and returns the index closest to the configured
target radius (default 5\,arcsec).  Results: Sinistro 1\,m new series
index~9 (27\,px $= 10.50$"), old series index~7 (27\,px $= 10.50$") ---
the \emph{same physical aperture} in both epochs, confirming the fix is
consistent.  Spectral/MuSCAT3 2\,m: new series index~11
(33\,px $= 10.03$"/8.91"); old series index~9 (same physical size).
\end{notebox}


\label{sec:aper}

Because the science aperture is not necessarily the largest, the script computes aperture
corrections: the median magnitude step between adjacent aperture indices is accumulated
outward to give a correction to infinite aperture for each of the 12 apertures.

The science aperture index is not fixed --- it depends on telescope family and epoch:
\begin{center}
\begin{tabular}{lll}
\toprule
Telescope & Condition & Aperture index \\
\midrule
1\,m & default & 7 \\
1\,m & MJD $> 66909$ & 9 \\
2\,m & MJD $\leq 60909$ & 10 \\
2\,m & MJD $> 60909$ & 12 \\
\bottomrule
\end{tabular}
\end{center}

\begin{notebox}
\textbf{Note.}  The physical aperture sizes in arcsec corresponding to each index depend on
the pixel scale.  For 1\,m Sinistro frames ($0.389$\,arcsec/px): index~7 corresponds to a
diameter of approximately 7\,arcsec; index~9 to approximately 9\,arcsec.  The 2026
expansion was introduced to maintain a consistent metric aperture as the LCO pixel scale
configuration changed. (\textit{Confirm exact MJD thresholds and arcsec values with
J.V.H.S.})
\end{notebox}

\subsection{Zeropoint fitting}

The zeropoint is estimated in two independent ways, both stored in the output:

\textbf{Bootstrap mean} (\texttt{zp0}, \texttt{zp0\_err}).  For each aperture, the median
shift (standard $-$ instrumental corrected) is bootstrap-resampled 50~times to give a
robust mean and uncertainty.

\textbf{\texttt{lmfit} fit} (\texttt{zp0\_fit}, \texttt{zp0\_err}).  A linear fit with
slope fixed to~1.0 and sigma-clipping of outliers yields a second independent estimate.

A minimum of 3~standard stars is required.  If fewer are available, or if the FWHM seeing
exceeds 9\,arcsec, the frame is flagged (\texttt{flag=1}) and stored with NaN photometry.

\subsection{Output columns in \texttt{lco\_latest.pkl}}

One row per FITS frame.  Array-valued columns (aperture vectors of length~12) are stored
as Python lists within DataFrame cells.

\begin{center}
\small
\begin{tabular}{lll}
\toprule
Column & Type & Description \\
\midrule
\texttt{flname}      & str    & Full path to BANZAI FITS file \\
\texttt{Object}      & str    & Target name from FITS header \\
\texttt{Filter}      & str    & LCO filter name \\
\texttt{MJD}         & float  & Mid-exposure MJD \\
\texttt{mag\_aper}   & list[12] & Aperture magnitudes (instrumental + offset) \\
\texttt{err\_aper}   & float  & Photometric error at science aperture \\
\texttt{zp\_mag}     & list[12] & Bootstrap zeropoints per aperture \\
\texttt{zp\_err}     & list[12] & Zeropoint uncertainties \\
\texttt{zp0\_fit}    & float  & \texttt{lmfit} zeropoint \\
\texttt{zp0\_err}    & float  & \texttt{lmfit} zeropoint uncertainty \\
\texttt{corrects}    & list[12] & Aperture corrections \\
\texttt{extra\_fair} & float  & Extinction + colour correction applied to AGN \\
\texttt{flag}        & int    & 0 = good; 1 = too few standards or bad seeing \\
\texttt{n\_std}      & int    & Number of standard stars used \\
\texttt{dist}        & float  & AGN--catalogue separation (arcsec) \\
\texttt{seeing}      & float  & Measured FWHM (arcsec) \\
\texttt{airmass}     & float  & Airmass $-1$ \\
\texttt{observatory} & str    & LCO site ID \\
\texttt{telid}       & str    & Telescope ID from filename \\
\texttt{skymag}      & float  & Sky background magnitude \\
\texttt{m\_fit}      & float  & Slope of linear photometric solution \\
\texttt{b\_fit}      & float  & Intercept of linear photometric solution \\
\texttt{X}, \texttt{Y} & float & AGN pixel position in image \\
\bottomrule
\end{tabular}
\end{center}

%%============================================================
\section{Stage 5 --- Plotting \& Web Output: \texttt{ava\_webplot.py}}
\label{sec:webplot}
%%============================================================

Reads \texttt{lco\_latest.pkl} and produces three sets of plots per target plus writes
calibrated light-curve data as text files for the website.

\subsection{Flux conversion}

Calibrated magnitudes are converted to flux density in mJy:
\begin{equation}
  f_\nu \,[\text{mJy}] = 10^{(m + \mathrm{ZP} - 8.9)\,/\,-2.5} \times 10^3
  \label{eq:flux}
\end{equation}
with uncertainties $\sigma_f = f_\nu \,/\, 1.0857 \times \sigma_m$.

\subsection{Quality cuts applied at plot time}

\begin{itemize}[noitemsep]
  \item \texttt{flag == 0} (calibration succeeded)
  \item \texttt{dist < 3.5}\,arcsec (AGN position match)
  \item \texttt{n\_std >= 3} (at least 3~standard stars)
\end{itemize}

\subsection{Plots produced}

\begin{enumerate}
  \item \textbf{Full light curve} (\texttt{lco\_latest\_\textit{obj}.pdf/png}).  One panel
        per filter over the full time baseline.  FLOYDS epoch times are marked as vertical
        ticks on the top panel.
  \item \textbf{Last 60~days} (\texttt{lco\_latest\_\textit{obj}\_last6.pdf/png}).  Same
        layout restricted to the 60~days preceding the most recent observation.  Y-axis
        limits are set from the 4th--96th flux percentile with 50\,\% padding.
  \item \textbf{Swift} (\texttt{lco\_latest\_\textit{obj}\_swift.pdf/png}).  Reads
        \texttt{ava\_swift\_updated.csv} and plots UV and X-ray light curves.
\end{enumerate}

\subsection{Diagnostic plots (\texttt{-d} flag)}

\begin{itemize}[noitemsep]
  \item Seeing, airmass, and zeropoint (transparency) vs time per site.
  \item Standard-star magnitude vs photometric RMS, with the AGN overlaid.  This is the
        primary internal photometric quality diagnostic.
\end{itemize}

%%============================================================
\section{Stage 6 --- Inter-calibration: \texttt{ava\_merge.py}}
\label{sec:merge}
%%============================================================

Photometry from different telescopes, cameras, and sites cannot be concatenated directly:
each has its own zeropoint offset.  \texttt{ava\_merge.py} uses
\texttt{pycali\_jvhs} --- a custom Python wrapper around the CALI MCMC
inter-calibration algorithm --- to solve for per-telescope flux scale factors and merge
the light curves into a single consistent time series.

\subsection{Input}

The script reads \texttt{<obj>\_edited.csv} from the parent directory.  This is the
calibrated light-curve file (essentially \texttt{lco\_latest.pkl} converted to CSV and
optionally hand-edited to remove bad epochs).  It must exist before merging can proceed.

\subsection{CALI algorithm}

For each filter, unique telescope identifiers (\texttt{telid}) are found.
\texttt{pycali\_jvhs.pycali()} is run with an MCMC chain of $10{,}000$ steps (extended
by $10{,}000$ on each retry if convergence is not achieved).  CALI models the light curve
as a common signal observed by each telescope with an additive offset and a flux scaling
factor, marginalising over these nuisance parameters.

\subsection{Output}

Merged light-curve files and plots are written to the \texttt{ava\_merging/} subdirectory,
then copied to \texttt{ava\_plots/} and \texttt{scp}\textquotesingle d to the web server.

\begin{warnbox}
\textbf{Migration target.}  \texttt{pycali\_jvhs} is a private package not on PyPI and
must be installed separately.  AVA~v2 should document the install procedure or replace
it with PyROA or CREAM.
\end{warnbox}

%%============================================================
\section{Data Products \& File Conventions}
\label{sec:products}
%%============================================================

All intermediate and final products are stored per-target under a fixed directory tree
rooted at \texttt{master\_dir} (currently \texttt{/data3/jvhs1/agn/}).

\begin{longtable}{>{\ttfamily}p{5.5cm}p{8.5cm}}
\toprule
File / folder & Contents \\
\midrule
apass\_\textit{obj}.csv & APASS reference star catalogue \\
apass\_pytics.csv & APASS catalogue reformatted for pytics column names \\
lco\_latest.pkl & Running pandas DataFrame of all calibrated photometry \\
lco\_latest\_stan.pkl & Running DataFrame of all calibrated standard stars \\
lco\_latest\_\textit{obj}.pdf/png & Full light-curve plot \\
lco\_latest\_\textit{obj}\_last6.pdf/png & Last-60-days light-curve plot \\
lco\_latest\_\textit{obj}\_swift.pdf/png & Swift UV + X-ray plot \\
lco\_latest\_\textit{obj}\_\textit{filt}\_sky.pdf/png & Diagnostic plots \\
lco\_latest\_\textit{obj}\_stan.pdf/png & Standard-star RMS plot \\
sextractor\_ava/ & SExtractor catalogues (\texttt{*\_cat.fits}) \\
ava\_merging/ & CALI workspace and merged output files \\
UVOT/ & Swift/UVOT data placeholder \\
default.sex, etc. & SExtractor configuration files \\
filter\_dict.txt & Filter name / wavelength / colour mapping \\
lco\_ext\_coefficients.csv & Extinction and colour-term coefficients \\
\bottomrule
\caption{Per-target working directory contents.}
\end{longtable}

\textbf{Raw FITS data} are stored at
\texttt{<master\_dir>/raw\_data/<obj>/YYYYMMDD/*e91.fits.fz}, one sub-folder per
observing night.

\begin{warnbox}
\textbf{Pickle portability.}  \texttt{lco\_latest.pkl} stores 12-element aperture
vectors as Python lists inside a pandas DataFrame.  Reading them requires
\texttt{np.stack(lco.mag\_aper[mask], axis=0)}.  This format is fragile across
pandas/numpy versions.  AVA~v2 should migrate to HDF5 or Parquet, or store photometry
directly in the \texttt{lcodb} SQLite database.
\end{warnbox}

%%============================================================
\section{Known Issues, Limitations \& Migration Notes}
\label{sec:issues}
%%============================================================

This section catalogues the technical debt that must be resolved before the pipeline can
be safely migrated to a new machine or handed to a new user.

\begin{enumerate}

\item \textbf{Hardcoded absolute paths.}  Every script contains literal paths to
      \texttt{/data/jvhs1/agn/} or \texttt{/data3/jvhs1/agn/}.  These are the single
      largest portability blocker.  AVA~v2 must replace all of them with paths read from
      \texttt{config/ava\_config.yaml}.

\item \textbf{APASS via web scrape.}  \texttt{ava\_initialise.py} fetches the APASS
      catalogue using the \texttt{mechanize} browser library.  Replacement:
      \texttt{astroquery.vizier.query\_region()} on APASS\,DR9 or Pan-STARRS\,PS1.

\item \textbf{Pickle-based state.}  \texttt{lco\_latest.pkl} with nested arrays is not
      portable across versions.  Replace with HDF5 / Parquet or the \texttt{lcodb}
      SQLite database.

\item \textbf{SExtractor as external binary.}  \texttt{os.system(`sex ...')} with no
      version checking.  Replace with \texttt{sep} (Python SExtractor) or
      \texttt{photutils}.

\item \textbf{Aperture index by magic number.}  Bare integer literals with inline MJD
      comparisons.  Must become a named config parameter with documented physical meaning.

\item \textbf{Credential in source code.}  \texttt{ava\_update.py} contains a hardcoded
      LCO archive username and password.  Replace with the token-based auth used by
      \texttt{lcodb}.

\item \textbf{SCP for delivery.}  Hardcoded hostname and path for plot delivery.
      Decouple pipeline from delivery mechanism.

\item \textbf{pycali\_jvhs undocumented dependency.}  Private package not on PyPI.
      Must be packaged or replaced with PyROA / CREAM.

\item \textbf{No test suite.}  No automated tests for any stage.  AVA~v2 should add
      unit tests for calibration arithmetic and integration tests on reference frames.

\item \textbf{Parallel execution via shell.}  Per-target parallelism via
      \texttt{nohup}~scripts provides no job management, logging, or failure recovery.
      AVA~v2 should use Python \texttt{multiprocessing} or a lightweight task queue.

\end{enumerate}

%%============================================================

%%============================================================
\section{Installation Guide}
\label{sec:install}
%%============================================================

This chapter documents how to install the AVA pipeline on a new machine from
scratch.  Follow the steps in order.  Where a command differs between macOS
and Linux, both variants are shown.

\subsection{Prerequisites}

The pipeline requires the following to be present on the machine before
installation begins:

\begin{itemize}
  \item \textbf{Conda} (Miniconda or Anaconda) --- used to create an isolated
        Python environment.  Install from
        \url{https://docs.conda.io/en/latest/miniconda.html}.
  \item \textbf{Git} --- to clone the repositories.
  \item \textbf{SExtractor} --- the external source-extraction binary.
        On macOS: \texttt{brew install sextractor}.
        On Linux: \texttt{sudo apt install sextractor} or compile from source.
  \item \textbf{SSH access} to the LCO data storage machine (if working
        remotely).
\end{itemize}

\subsection{Step 1 --- Clone the repositories}

\begin{lstlisting}[style=shell]
# Choose a base directory (adjust to taste)
BASE=~/astro

mkdir -p $BASE && cd $BASE

# AVA main repository
git clone https://github.com/jvhs1/AVA.git

# banzai-floyds (spectroscopy pipeline)
# Use the rm-improvements branch for the RM-optimised version
git clone https://github.com/LCOGT/banzai-floyds.git floyds/banzai-floyds
cd $BASE/floyds/banzai-floyds
git checkout rm-improvements
cd $BASE
\end{lstlisting}

\subsection{Step 2 --- Create the conda environment}

\begin{lstlisting}[style=shell]
cd $BASE/AVA
conda env create -f environment.yml
conda activate ava
\end{lstlisting}

This creates an environment named \texttt{ava} with Python~3.12 and all
required packages.  To verify:
\begin{lstlisting}[style=shell]
python -c "import ava; print(ava.__version__)"
\end{lstlisting}

\subsection{Step 3 --- Install AVA and its sibling packages}

\begin{lstlisting}[style=shell]
# Install AVA itself in editable mode
pip install -e "$BASE/AVA"

# Install banzai-floyds in editable mode (spectroscopy only)
pip install -e "$BASE/floyds/banzai-floyds"

# pycali_jvhs (inter-calibration; install from local copy)
# pip install -e /path/to/pycali_jvhs
# Ask J.V.H.S. for the source if not already present.
\end{lstlisting}

After this step the \texttt{lcodb} command is available (it is bundled inside
AVA):
\begin{lstlisting}[style=shell]
lcodb --help
\end{lstlisting}

\subsection{Step 4 --- Set the LCO API token}

Obtain your personal API token from
\url{https://observe.lco.global} $\to$ Profile $\to$ API~token.

\begin{lstlisting}[style=shell]
echo 'YOUR_TOKEN_HERE' > ~/.lco_token
chmod 600 ~/.lco_token
\end{lstlisting}

\begin{warnbox}
Never put the token directly in a config file or commit it to git.
\end{warnbox}

\subsection{Step 5 --- Create and configure your local config file}

\begin{lstlisting}[style=shell]
cd $BASE/AVA
cp config/ava_config.yaml config/ava_config.local.yaml
\end{lstlisting}

Open \texttt{config/ava\_config.local.yaml} in an editor and update the
\texttt{paths} section to match the new machine.  The keys that \emph{must}
be set are listed in Table~\ref{tab:paths}.

\begin{table}[h!]
\centering
\small
\begin{tabular}{>{\ttfamily}p{5.2cm} p{4.5cm} p{4.0cm}}
\toprule
Key & Description & Production value \\
\midrule
paths.photometry\_root   & Root of photometry data tree & \texttt{/data3/jvhs1/agn} \\
paths.photometry\_raw    & Raw frames subfolder         & \texttt{raw\_data} \\
paths.spectra\_root      & Root of spectroscopy data tree & \texttt{/data2/jvhs1/agn} \\
paths.spectra\_raw       & Raw FLOYDS frames subfolder  & \texttt{raw\_data} \\
paths.plots\_dir         & Shared plot output directory & \texttt{/data3/jvhs1/agn/ava\_plots} \\
paths.swift\_db          & Swift/UVOT merged CSV        & \texttt{/data3/jvhs1/agn/ava\_swift\_updated.csv} \\
paths.template\_dir      & SExtractor config templates  & \texttt{/data3/jvhs1/agn/template\_pipeline} \\
paths.floyds\_pipeline   & banzai-floyds working copy   & \texttt{\$BASE/floyds/banzai-floyds} \\
paths.database           & AVA SQLite master DB         & \texttt{data/ava.db} \\
paths.ava\_csv           & AVA master CSV               & \texttt{data/ava/ava.csv} \\
\bottomrule
\end{tabular}
\caption{Path keys that must be set for each machine.  Relative paths are
  resolved from \texttt{ava\_root} (auto-detected as the repository root if
  left blank).}
\label{tab:paths}
\end{table}

Point AVA at the local config:
\begin{lstlisting}[style=shell]
# Option A: environment variable (add to ~/.bashrc or ~/.zshrc)
export AVA_CONFIG="$BASE/AVA/config/ava_config.local.yaml"

# Option B: symlink the local copy over the default
ln -sf ava_config.local.yaml config/ava_config.yaml
\end{lstlisting}

Verify that the config loads and all paths resolve:
\begin{lstlisting}[style=python]
from ava.config import cfg
print(cfg)
# AVAConfig(
#   config  = /home/.../config/ava_config.local.yaml
#   root    = /home/.../AVA
#   phot    = /data3/jvhs1/agn
#   spec    = /data2/jvhs1/agn
#   db      = /home/.../AVA/data/ava.db
# )
\end{lstlisting}

\subsection{Step 6 --- Initialise the LCO archive databases}

This step builds the \texttt{lcodb} archive databases by crawling the LCO
archive API.  It is a one-time operation that takes between a few minutes
(spectra) and several hours (photometry headers) depending on the data volume.

\begin{lstlisting}[style=shell]
# Crawl all KEY proposals (metadata only; no FITS files downloaded)
lcodb sync

# Download FITS headers for all spectra (~10^4 frames, minutes)
lcodb headers --db spectra

# Compute airmass for frames without it in the header
lcodb airmass

# Build science->standard and spectrum->arc/flat matching tables
lcodb match

# Link LCO target names to AVA sources
lcodb link

# Check status
lcodb status
\end{lstlisting}

After this step, the archive databases are ready at the paths set in
\texttt{config/lcodb\_config.yaml} (default: \texttt{data/lcodb/}).

\subsection{Step 7 --- Migrate the existing AVA database}

If migrating from an existing machine, copy the master SQLite database and
run the migration script to import it into the new schema:

\begin{lstlisting}[style=shell]
# Copy from the old machine (adjust path)
scp jvhs1@oldmachine:/path/to/ava.sqlite data/ava/ava.sqlite.old

# Migrate into the new AVA database
python -m ava.database.migrate \
    --source data/ava/ava.sqlite.old \
    --dest data/ava.db
\end{lstlisting}

\subsection{Step 8 --- Schedule the daily update}

The daily update fetches new frames from the LCO archive and updates the
databases.  Add the following to \texttt{crontab -e}:

\begin{lstlisting}[style=shell]
# Run every day at 09:30 UTC
30 9 * * * /path/to/AVA/scripts/lcodb_daily.sh
\end{lstlisting}

Or on macOS using \texttt{launchd}, create a plist pointing to the same script.

\subsection{Step 9 --- Verify the installation}

\begin{lstlisting}[style=python]
from ava.archive import get_archive
a = get_archive()
print(a)              # Archive(spectra=ok, photometry=ok, ava=ok)
a.overview()          # Should show the full AGN table
a.summary('Fairall 9')
\end{lstlisting}

Run the test suite to confirm all units pass:
\begin{lstlisting}[style=shell]
pytest tests/ -v
\end{lstlisting}

\subsection{Data directory structure on the production machine}

For reference, the expected layout on the production machine
(\texttt{/data2}, \texttt{/data3}) is:

\begin{lstlisting}[style=shell]
# Photometry  (--  /data3/jvhs1/agn/ --)
/data3/jvhs1/agn/
  raw_data/
    <target_slug>/
      YYYYMMDD/
        *e91.fits.fz        # BANZAI-reduced frames
  <target_slug>/
    apass_<target>.csv      # APASS reference catalogue
    apass_pytics.csv        # Reformatted for pytics
    lco_latest.pkl          # Calibrated photometry (all epochs)
    lco_latest_stan.pkl     # Calibrated standard stars
    lco_latest_*.pdf/png    # Light-curve plots
    sextractor_ava/
      *_cat.fits            # SExtractor catalogues
    ava_merging/            # CALI workspace
  ava_plots/                # Plots copied here before web delivery
  ava_swift_updated.csv     # Swift/UVOT merged database
  template_pipeline/        # SExtractor config templates

# Spectroscopy  (--  /data2/jvhs1/agn/ --)
/data2/jvhs1/agn/
  raw_data/
    <target_slug>/
      YYYYMMDD/
        *e91.fits.fz        # BANZAI FLOYDS frames
  <target_slug>/            # banzai-floyds reduction products
\end{lstlisting}


%%============================================================
\section{Spectroscopic Archive: Standards \& Coverage}
\label{sec:spec-archive}
%%============================================================

\textit{Last updated: 2026-11-15. Source:
\texttt{floyds/lco\_archive\_db/db/spectra.sqlite} (252\,MB).}

\subsection{Overall statistics}

\begin{table}[h!]
\centering
\small
\begin{tabular}{lr}
\toprule
Metric & Value \\
\midrule
Science spectra (raw, \texttt{rlevel=0}) & 9{,}524 \\
Standard spectra (raw) & 20{,}712 \\
Science spectra with $\geq 1$ matched standard & \textbf{9{,}524 (100\,\%)} \\
Science spectra with no standard & \textbf{0} \\
Total \texttt{std\_match} rows & 63{,}615 \\
BANZAI 1d products available & 3{,}110 / 21{,}477 (14.5\,\%) \\
Total FLOYDS exposure time & 1{,}926\,h across 8{,}860 frames, 5{,}139 epochs \\
Targets with any FLOYDS data & 57 of 71 \\
\bottomrule
\end{tabular}
\caption{FLOYDS spectroscopic archive summary (2026-11-15).}
\end{table}

\begin{notebox}
\textbf{Full standard coverage.}  Every one of the 9{,}524 science spectra has
at least one matched flux standard in \texttt{spectra.sqlite}.  No coverage
gaps exist in the current archive.
\end{notebox}

\subsection{Spectrograph and slit usage}

\begin{table}[h!]
\centering
\small
\begin{tabular}{llrll}
\toprule
Instrument & Site & Science spectra & First obs & Last obs \\
\midrule
en06 (FLOYDS-N) & ogg (Haleakala)     & 5{,}656 & 2014-08-05 & 2026-11-14 \\
en12 (FLOYDS-S) & coj (Siding Spring) & 2{,}570 & 2019-04-14 & 2026-11-11 \\
en05 (FLOYDS-S, retired) & coj        &   634   & 2014-05-19 & 2019-02-03 \\
\bottomrule
\end{tabular}
\caption{Science spectra by spectrograph.}
\end{table}

Slit-width distribution: 5{,}854 at 6\,arcsec; 1{,}380 at 2\,arcsec;
126 at 1.6\,arcsec; 10 at 1.2\,arcsec; 1{,}490 with no slit keyword in the
header.

\subsection{Standard selection}

The \texttt{spectrum\_standards} view in \texttt{spectra.sqlite} assigns one
standard per science spectrum using this priority order:

\begin{enumerate}
  \item \textbf{\texttt{same\_slit}} --- nearest standard taken with the same
        slit width (2{,}849 spectra, mean $\langle\Delta t\rangle = 176$\,h).
  \item \textbf{\texttt{nearest\_2arcsec}} --- nearest 2\,arcsec standard when
        no same-slit option exists (4{,}280 spectra, $\langle\Delta t\rangle =
        48$\,h).
  \item \textbf{\texttt{nearest\_any}} --- nearest standard of any slit
        (1{,}731 spectra, $\langle\Delta t\rangle = 30$\,h).
\end{enumerate}

Standards are \emph{never taken in the same observing block} as the science
(\texttt{same\_block = 0} throughout); FLOYDS standards are submitted as
separate queue requests.  The minimum $\Delta t$ is typically $\approx
12$\,minutes.

\subsection{Targets with large average standard $\Delta t$}
\label{sec:std-dt}

The targets below had sparse standard coverage during their campaigns and
may carry systematic flux calibration offsets.

\begin{table}[h!]
\centering
\small
\begin{tabular}{lrr}
\toprule
Target & Science spectra & $\langle|\Delta t|\rangle$ (hours) \\
\midrule
Mrk\,1220      &  86 & 552 \\
3C\,120        & 153 & 299 \\
Mrk\,110       & 217 & 232 \\
SDSS\,J2353    &  13 & 168 \\
PKS\,2004-447  &   7 &  88 \\
PG\,1302-102   &  99 &  71 \\
NGC\,3783      & 177 &  58 \\
Mrk\,1239      &  49 &  55 \\
Mrk\,766       &  50 &  50 \\
\bottomrule
\end{tabular}
\caption{Targets with $\langle|\Delta t|\rangle > 50$\,h to nearest standard.
  These are predominantly old campaigns (KEY2014A-002, KEY2018B-001) where
  dedicated FLOYDS standards were taken less frequently.}
\end{table}

\subsection{Targets with no FLOYDS data (14)}

SDSS-RM201, SDSS-RM562, PG\,1310-108, 2MASS\,J21368-6438,
2MASX\,J21075970-3754087, PG\,2304+043, RX\,J22198-0935, NGC\,4395,
PKS\,0558-504, NGC\,2617, Zw\,535-12, SDSS\,J1430+2303, Mrk\,704,
IGR\,J17062-6143.

\subsection{Database update log}

\begin{table}[h!]
\centering
\small
\begin{tabular}{lll}
\toprule
File & Location & Updated \\
\midrule
\texttt{ava.db}                 & \texttt{AVA/data/} & 2026-11-15 \\
\texttt{ava.csv}                & \texttt{AVA/data/ava/} & 2026-11-15 \\
\texttt{ava\_link\_overrides.csv} & \texttt{AVA/data/ava/} & 2026-11-15 \\
\texttt{spectra.sqlite}         & \texttt{floyds/lco\_archive\_db/db/} & 2026-11-15 \\
\texttt{photometry.sqlite}      & \texttt{floyds/lco\_archive\_db/db/} & live (daily) \\
\bottomrule
\end{tabular}
\caption{AVA data files and their last update dates.  \texttt{spectra.sqlite}
  and \texttt{photometry.sqlite} remain in the \texttt{floyds/} sandbox and
  are accessed read-only via \texttt{get\_archive()}.  A snapshot of
  \texttt{spectra.sqlite} taken before the 2026-11-15 calibration update is
  kept as \texttt{spectra\_before\_calib\_20261115.sqlite}.}
\end{table}


\appendix
\section{Running the Pipeline}
\label{app:run}
%%============================================================

\textbf{Full run for a single target on the latest data:}
\begin{lstlisting}[style=shell]
python ava_update.py -d -e -c -p -list target_list.txt -project latest
\end{lstlisting}

\textbf{Monitor progress:}
\begin{lstlisting}[style=shell]
ps aux | grep 'python download_'
ps aux | grep 'python ava_sextractor'
ps aux | grep 'python ava_calibrate'
tail -100 lco_Fairall_9_logall.out
\end{lstlisting}

\textbf{Force full recalibration (e.g.\ after aperture change):}
\begin{lstlisting}[style=shell]
python ava_calibrate.py -obj Fairall_9 -s -j   # -j applies 2026 aperture fix
\end{lstlisting}

\textbf{Inter-calibration:}
\begin{lstlisting}[style=shell]
python ava_update.py -m -s -list target_list.txt
\end{lstlisting}

%%============================================================
\section{Glossary}
\label{app:glossary}
%%============================================================

\begin{description}
  \item[APASS] AAVSO Photometric All-Sky Survey.  Provides $BVgri$ magnitudes for
        stars used as photometric standards.
  \item[BANZAI] LCO automated image-reduction pipeline.  Delivers bias-subtracted,
        flat-fielded, astrometry-solved FITS frames at \texttt{rlevel=91}.
  \item[CALI] Cross-calibration algorithm for multi-telescope light curves.  Fits
        per-telescope flux scale factors using MCMC.
  \item[MJD] Modified Julian Date. MJD = JD $-$ 2400000.5.
  \item[ZP] Photometric zeropoint.  The magnitude of a source that would produce
        1 ADU/s in the instrument.
  \item[SExtractor] Source-Extractor; software for source detection and aperture
        photometry in astronomical images.
\end{description}

\vfill
\hrule
\vspace{4pt}
{\small
Created: 2026-11-15 \quad
Author: J.V.\ Hernandez Santisteban \quad
Pipeline version: v1.3 (legacy) / AVA~v2 (in development)
}

\end{document}
